\section*{题目审阅记录}
\phantomsection\addcontentsline{toc}{section}{题目审阅记录}
以下为全书精校发现的问题及核查依据，不代替全书逐题完整解答。正文保留原题，修订建议未自动改入题面。每条卷号和题块均可点击返回正文。
\Needspace{8\baselineskip}
\subsection*{R001\quad 二.1 原书保留了未解析的公式引用}
\noindent\auditref{p036-b23}{第 32 套试卷；原书第 35 页}\quad 题面异常\par
扫描页确有以两个问号代替的公式引用。这是原书题面原有的占位符，不是本稿编译产生的未定义引用；按上下文很可能指前一行的方程。
\par\textbf{修订建议：}保留原题并记录；修订版可明确写为“由上述方程确定”，不凭空添加公式编号。\par
\Needspace{8\baselineskip}
\subsection*{R002\quad 二.2 原书同样使用问号引用}
\noindent\auditref{p036-b24}{第 32 套试卷；原书第 35 页}\quad 题面异常\par
扫描页确有问号引用，所指很可能是上方 $e^{x+2y+3z}+xyz=1$，但题面没有有效的公式编号。
\par\textbf{修订建议：}保留原文；若制作修订版，可明确关联上述方程。\par
\Needspace{8\baselineskip}
\subsection*{R003\quad 第9题原书保留了疑问括注}
\noindent\auditref{p101-b13}{第 100 套试卷；原书第 100 页}\quad 题面异常\par
原题把偏导、可微和连续的要求混写，并附三个问号。计算原点方向导数的要求明确，前半句需要明确各性质所针对的点和是否为讨论题。
\par\textbf{修订建议：}不将问号当成重排新增的引用错误；在修订版中明确设问。\par
\Needspace{8\baselineskip}
\subsection*{R004\quad 二.2 的可导结论范围过大}
\noindent\auditref{p012-b18}{第 4 套试卷；原书第 11 页}\quad 条件缺失\par
由题设只能得到 $f(0)=0$ 和 $f'(0)=3/4$。取 $f(x)=3x/4+x^2|x-1|$，题设极限成立，但在 $x=1$ 不可导。原文“求”亦疑为“证明”。
\par\textbf{修订建议：}将结论明确为“在 $x=0$ 处可导”；是否要求全区间可导需查原始回忆。\par
\Needspace{8\baselineskip}
\subsection*{R005\quad 八题的反常积分没有实参数收敛域}
\noindent\auditref{p013-b00}{第 4 套试卷；原书第 12 页}\quad 结论错误\par
$x\to0$ 要求 $p<1$；无穷端因 $\cos^2x$ 的正平均值要求 $p>1$，两者不相容。因此不能在所述区间讨论连续性。
\par\textbf{修订建议：}可能应为 $\sin^2x/x^p$，但这只是推测，正文保留 $\cos^2x$。\par
\Needspace{8\baselineskip}
\subsection*{R006\quad 二题仅有连续性，不能推出驻点}
\noindent\auditref{p013-b06}{第 5 套试卷；原书第 12 页}\quad 条件缺失\par
取 $f(x)=3/2-x$（$0\le x\le1$），$f(x)=(x-1)/2+1/2$（$1\le x\le2$）。满足 $f(0)+f(1)=2,f(2)=1$，但可导点的导数均不为零，$x=1$ 不可导。
\par\textbf{修订建议：}补充在 $(0,2)$ 内可导。\par
\Needspace{8\baselineskip}
\subsection*{R007\quad 第6题不存在所求有限上界常数}
\noindent\auditref{p014-b10}{第 6 套试卷；原书第 13 页}\quad 结论错误\par
令 $x=y=t\to+\infty$，上界要求 $B\ge2t^2\ln(2t)$，不存在有限的 $B$。
\par\textbf{修订建议：}核实变量范围、对数自变量或不等号方向；不能自行补数据。\par
\Needspace{8\baselineskip}
\subsection*{R008\quad 十五(1) 的序列范围与无界方向}
\noindent\auditref{p021-b20}{第 15 套试卷；原书第 20 页}\quad 条件缺失\par
函数仅定义在 $[a,b]$，不能直接选 $x_n\in[0,1]$。若“$\infty$”指 $+\infty$，无界也未必向上无界，例如 $f(x)=-1/x$（$x>0$）且 $f(0)=0$。
\par\textbf{修订建议：}改为 $x_n\in[a,b]$ 且 $|f(x_n)|\to+\infty$；若原意为向上无界需明确。\par
\Needspace{8\baselineskip}
\subsection*{R009\quad 十五(2) 无界聚集点不一定在内部}
\noindent\auditref{p021-b21}{第 15 套试卷；原书第 20 页}\quad 结论错误\par
在 $[0,1]$ 上令 $f(0)=0,f(x)=1/x$（$x>0$）。无界只聚集于端点 $0$，任何内部点都有一个使函数有界的小邻域。
\par\textbf{修订建议：}将 $c\in(a,b)$ 改为 $c\in[a,b]$。\par
\Needspace{8\baselineskip}
\subsection*{R010\quad 第7题要求的 Fourier 逐点极限缺少条件}
\noindent\auditref{p023-b17}{第 17 套试卷；原书第 22 页}\quad 条件缺失\par
可积性足以确定最小二乘系数，但不能保证 Fourier 部分和在每点收敛；一般连续函数的 Fourier 级数也可能在某点发散。
\par\textbf{修订建议：}说明只求均方极限，或增加使所求点处收敛的正则性条件；两种修订含义不同。\par
\Needspace{8\baselineskip}
\subsection*{R011\quad 第8题未排除恒零迭代}
\noindent\auditref{p023-b20}{第 17 套试卷；原书第 22 页}\quad 条件缺失\par
取 $f(x)=x-x^3/6$，$a_n\equiv0$ 满足全部条件，此时 $na_n^2\to0$；非零且趋零的轨道则由三阶展开给出 $na_n^2\to3$。
\par\textbf{修订建议：}如要求唯一数值 $3$，需排除某项等于零的轨道。\par
\Needspace{8\baselineskip}
\subsection*{R012\quad 二题一致连续性结论相反}
\noindent\auditref{p032-b13}{第 27 套试卷；原书第 31 页}\quad 结论错误\par
取 $x_n=(2\pi n)^{-1}$、$y_n=((2n+1)\pi)^{-1}$。两列距离趋零，而函数值之差趋于 $2$，故不一致连续。
\par\textbf{修订建议：}将“证明一致连续”改为“证明不一致连续”，仍需核原始回忆。\par
\Needspace{8\baselineskip}
\subsection*{R013\quad 六题函数列不一致收敛}
\noindent\auditref{p032-b18}{第 27 套试卷；原书第 31 页}\quad 结论错误\par
逐点极限为 $0$，但在 $x=1/(n+1)$ 处有 $nx(1-x)^n=(n/(n+1))^{n+1}\to e^{-1}$，因此上确界不趋零。
\par\textbf{修订建议：}改为讨论一致收敛性，或将结论改为非一致收敛。\par
\Needspace{8\baselineskip}
\subsection*{R014\quad 十题函数极限不能推出所写导数极限}
\noindent\auditref{p037-b09}{第 32 套试卷；原书第 36 页}\quad 结论错误\par
取 $f(x)=1/x$，则 $f(x)\to+\infty$，但 $f'(x)=-1/x^2\to-\infty$。即使改成负号，一般情形也不能仅凭函数极限推出导数极限。
\par\textbf{修订建议：}重新核实所求结论，不仅是机械改正负号。\par
\Needspace{8\baselineskip}
\subsection*{R015\quad 第14题两个 L1 条件不保证乘积可积}
\noindent\auditref{p039-b20}{第 35 套试卷；原书第 38 页}\quad 条件缺失\par
取 $[a,b]=[0,1]$，$f(x)=x^{-2/3}$（原点任定有限值），$g$ 为 $x^{-2/3}$ 在 $(0,1)$ 的 1 周期延拓。二者可积，但对每个正整数 $n$，$f(x)g(nx)$ 在原点如常数乘 $x^{-4/3}$，积分发散。
\par\textbf{修订建议：}可增加其中一个函数有界等足够条件；不能把现题当作一般 L1 引理。\par
\Needspace{8\baselineskip}
\subsection*{R016\quad 五题矩阵与向量的收缩表达式异常}
\noindent\auditref{p046-b24}{第 42 套试卷；原书第 45 页}\quad 题面异常\par
最后一项原图确印为 $\nabla u\cdot(D^2u\cdot\Delta u)$。$\Delta u$ 是标量，此式的结果类型与左侧标量不符；通常对应的是 $\nabla u\cdot(D^2u\,\nabla u)$。
\par\textbf{修订建议：}将最后一个 $\Delta u$ 列作疑似 $\nabla u$ 笔误，等待确认；正文保留原式。\par
\Needspace{8\baselineskip}
\subsection*{R017\quad 第2题参数范围未给出}
\noindent\auditref{p050-b08}{第 47 套试卷；原书第 49 页}\quad 条件缺失\par
实对数要求比值为正。若 $|a|>1$，积分区间中会出现比值为负的区段；$a=\pm1$ 又有端点问题。
\par\textbf{修订建议：}明确实参数范围，并说明是否按广义积分解释端点。\par
\Needspace{8\baselineskip}
\subsection*{R018\quad 三题导数积分缺少可导假设}
\noindent\auditref{p051-b08}{第 48 套试卷；原书第 50 页}\quad 条件缺失\par
原书仅写连续并另注“应该是连续可微”。一般连续函数的 $f'$ 不能据此定义，题中的两个积分也无保证。
\par\textbf{修订建议：}建议补连续可微；保留原编括注。\par
\Needspace{8\baselineskip}
\subsection*{R019\quad 二.5 的严格不等号及开区间条件}
\noindent\auditref{p053-b07}{第 50 套试卷；原书第 52 页}\quad 题面异常\par
原式允许 $x_1=x_2$，此时严格不等号变为 $0<0$，题设本身矛盾。即便补 $x_1\ne x_2$，局部 Lipschitz 也不能推出开区间上的统一常数，例如 $(0,1)$ 上的 $1/x$。
\par\textbf{修订建议：}同时核实严格号、不同点限制与区间是否应为闭区间。\par
\Needspace{8\baselineskip}
\subsection*{R020\quad 八题在分段连接点不连续}
\noindent\auditref{p055-b21}{第 53 套试卷；原书第 54 页}\quad 结论错误\par
左端在 $x=\pi/2$ 取值 $(\pi/2)^\alpha\cos(2/\pi)\ne0$，右极限为 $0$，故不可能在整个 $(0,+\infty)$ 上一致连续。
\par\textbf{修订建议：}核实左段的余弦自变量及分段点。\par
\Needspace{8\baselineskip}
\subsection*{R021\quad 四(2) Green 表示公式第二项符号}
\noindent\auditref{p063-b04}{第 62 套试卷；原书第 62 页}\quad 结论错误\par
原书定义 $\boldsymbol r=\boldsymbol x-\boldsymbol x_0$，外法向为 $\boldsymbol n$。由 Green 第二恒等式得 $4\pi f(\boldsymbol x_0)=\int_S[f\cos(\boldsymbol r,\boldsymbol n)/r^2-(1/r)\partial_n f],dS$，原题写加号。
\par\textbf{修订建议：}第二项加号应为减号。\par
\Needspace{8\baselineskip}
\subsection*{R022\quad 四(3) 积分恒等式符号相反}
\noindent\auditref{p063-b08}{第 62 套试卷；原书第 62 页}\quad 结论错误\par
取 $f=(x^2+y^2+z^2)/6$、$g\equiv1$。左侧为 $4\pi/15$，原右侧为 $\pi(2/5-2/3)=-4\pi/15$。
\par\textbf{修订建议：}右侧 $(t-1)$ 应改为 $(1-t)$。\par
\Needspace{8\baselineskip}
\subsection*{R023\quad 四(2)(3) 估计对象疑漏导数撇号}
\noindent\auditref{p081-b00}{第 81 套试卷；原书第 80 页}\quad 结论错误\par
取 $f\equiv1$，可取 $M_0=1,M_2=0$，则(2)在 $h>1$ 失败，(3)变为 $1\le0$。原图的两处均印作 $|f(x)|$。
\par\textbf{修订建议：}通常应估计 $|f\,\prime(x)|$；两处一并核实。\par
\Needspace{8\baselineskip}
\subsection*{R024\quad 第7(1)题 Chebyshev 不等号反向}
\noindent\auditref{p101-b08}{第 100 套试卷；原书第 100 页}\quad 结论错误\par
取 $[a,b]=[0,1]$、$f(x)=g(x)=x$，满足同单调条件，原结论变为 $1/4\ge1/3$。
\par\textbf{修订建议：}不等号应为 $\le$。\par
\Needspace{8\baselineskip}
\subsection*{R025\quad 第4题极限结论缺少导数或其他条件}
\noindent\auditref{p112-b19}{第 114 套试卷；原书第 111 页}\quad 结论错误\par
取 $f\equiv1$，则 $x^{2025}f'(x)=0$ 单调，然而 $(x\ln x)f(x)\to+\infty$。
\par\textbf{修订建议：}核实结论中的 $f$ 是否应为 $f\,\prime$；不能仅添加“单调严格”掩盖问题。\par
\Needspace{8\baselineskip}
\subsection*{R026\quad 第7(2)题三维缩放系数少了2}
\noindent\auditref{p113-b03}{第 114 套试卷；原书第 112 页}\quad 结论错误\par
题目明确为三维。取调和函数 $u\equiv1$，则 $I(r)=0$，变量代换给出 $H(r)=C r^{2\alpha+3}$，故 $H'=(2\alpha+3)H/r$，不是 $(2\alpha+1)H/r$。
\par\textbf{修订建议：}将系数 $2\alpha+1$ 改为 $2\alpha+3$。\par
\Needspace{8\baselineskip}
\subsection*{R027\quad 第3(1)题 Cesaro 逆命题为假}
\noindent\auditref{p121-b13}{第 126 套试卷；原书第 120 页}\quad 题面异常\par
取 $a_n=(-1)^n$，算术平均趋于 $0$，原数列不收敛。本题位于“解答如下问题”下，若原意为判断题则这是应答反例，不应改题；若要求证明则缺条件。
\par\textbf{修订建议：}核实题型，或补单调等可使逆命题成立的条件。\par
\Needspace{8\baselineskip}
\subsection*{R028\quad 第6(1)题级数在全实轴上并未定义}
\noindent\auditref{p121-b18}{第 126 套试卷；原书第 120 页}\quad 结论错误\par
当 $n=1$ 时分母为 $x+1$，在 $x=-1$ 为零，故全实轴一致收敛的结论不成立。
\par\textbf{修订建议：}需重新确定定义域，并检查所有奇数项的极点；不能只删去 $-1$。\par
\Needspace{8\baselineskip}
\subsection*{R029\quad 第7(1)题平移后的幂积分公式错误}
\noindent\auditref{p137-b03}{第 142 套试卷；原书第 136 页}\quad 结论错误\par
令 $u=x-(a+b)/2$，正确结果为 $(b-a)^{n+1}/[2^n(n+1)]$。如 $a=1,b=2,n=1$，左侧为 $1/4$，原右侧为 $3/4$。
\par\textbf{修订建议：}将分子 $b^{n+1}-a^{n+1}$ 改为 $(b-a)^{n+1}$。\par
\Needspace{8\baselineskip}
\subsection*{R030\quad 三题级数下限包含除以零项}
\noindent\auditref{p138-b14}{第 144 套试卷；原书第 137 页}\quad 题面异常\par
原书确印 $\sum_{n=0}^{\infty}1/n^2$，其第零项无定义；积分是有限的 $-\pi^2/6$。
\par\textbf{修订建议：}下限应为 $n=1$。\par
\Needspace{8\baselineskip}
\subsection*{R031\quad 八题无穷积分的零端未受控制}
\noindent\auditref{p139-b03}{第 144 套试卷；原书第 138 页}\quad 条件缺失\par
取 $f(x)=1/x^2$，在 $(0,+\infty)$ 连续、单调且趋零，但 $f(x)\sin(nx)\sim n/x$，零端积分发散。
\par\textbf{修订建议：}增加足以控制零端的条件，例如 $f$ 在零端有界；再核实所需极限。\par
\Needspace{8\baselineskip}
\subsection*{R032\quad 填空(9) 交线不是所写闭曲线}
\noindent\auditref{p147-b00}{第 153 套试卷；原书第 146 页}\quad 题面异常\par
柱面 $x^2+y^2=R^2$ 与不含 $z$ 的平面 $x/R+y/H=1$ 相交为沿 $z$ 方向的母线，并不构成闭合椭圆；闭路积分和指定方向无从成立。
\par\textbf{修订建议：}核实平面方程是否有 $z$ 项，不猜测具体分母。\par
\Needspace{8\baselineskip}
\subsection*{R033\quad 二题正切的周期性允许跳跃}
\noindent\auditref{p148-b06}{第 154 套试卷；原书第 147 页}\quad 结论错误\par
令 $h(x)=1-x/10$，取 $f(x)=h(x)$（$x<1/2$），$f(x)=h(x)+\pi$（$x\ge1/2$）。$\tan f=\tan h$ 严格递减；各段上 $(f\tan x)'=[f-(\sin x\cos x)/10]/\cos^2x>0$，分段点又向上跳跃，故 $f\tan x$ 严格递增，但 $f$ 不连续。
\par\textbf{修订建议：}需补 $f$ 的值域位于正切同一连续分支等限制，或核对题面两个表达式是否识别无误。\par
\Needspace{8\baselineskip}
\subsection*{R034\quad 四题振荡函数未必有趋零子列}
\noindent\auditref{p158-b12}{第 165 套试卷；原书第 157 页}\quad 结论错误\par
取 $f(x)=2+\sin(1/x)$。有界可导，零端极限不存在，所有函数值均在 $[1,3]$，不可能存在趋零函数值子列。
\par\textbf{修订建议：}核实所求是 $f\,\prime(x_n)\to0$，还是某个聚点值；两者不能混写。\par
\Needspace{8\baselineskip}
\subsection*{R035\quad 第3题缺少端点联系}
\noindent\auditref{p163-b03}{第 170 套试卷；原书第 162 页}\quad 条件缺失\par
取 $f\equiv1$、$\lambda=1$，则 $f'+\lambda f\equiv1$，没有零点。
\par\textbf{修订建议：}通常需补 $f(a)=f(b)=0$ 等使加权 Rolle 定理适用的条件。\par
\Needspace{8\baselineskip}
\subsection*{R036\quad 第8题连续且可积不能推出 xf趋零}
\noindent\auditref{p163-b08}{第 170 套试卷；原书第 162 页}\quad 结论错误\par
令 $x_n=2^{-n}$，在 $x_n$ 处叠加高度 $1/x_n$、半宽 $x_n^3$ 的不相交三角尖峰，并加上 $x^{-1/2}$。函数连续、零端趋 $+\infty$ 且可积，但 $x_nf(x_n)\ge1$。
\par\textbf{修订建议：}例如增加单调性可修复；需核原题条件。\par
\Needspace{8\baselineskip}
\subsection*{R037\quad 第12题无界点可能只在端点}
\noindent\auditref{p163-b13}{第 170 套试卷；原书第 162 页}\quad 结论错误\par
在 $[0,1]$ 上取 $f(0)=0,f(x)=1/x$（$x>0$）。任意内部点都有函数有界的小邻域。
\par\textbf{修订建议：}$x_0\in(a,b)$ 应为 $x_0\in[a,b]$。\par
\Needspace{8\baselineskip}
\subsection*{R038\quad 第12题对称割线没有检测中心函数值}
\noindent\auditref{p170-b16}{第 179 套试卷；原书第 169 页}\quad 条件缺失\par
取 $f(x_0)=1$，$x\ne x_0$ 时 $f(x)=0$。题中二元极限为 $0$，但 $f$ 在 $x_0$ 不连续，更不可导。
\par\textbf{修订建议：}补充 $f$ 在 $x_0$ 连续。\par
\Needspace{8\baselineskip}
\subsection*{R039\quad 四题零解结论缺少初值}
\noindent\auditref{p186-b02}{第 195 套试卷；原书第 185 页}\quad 条件缺失\par
取 $f\equiv1,g\equiv0$，则 $|f'+fg|=0\le|f|$，但 $f$ 不恒零。
\par\textbf{修订建议：}至少需类似 $f(0)=0$ 的初值，并明确 $f$ 的可导性。\par
\Needspace{8\baselineskip}
\subsection*{R040\quad 六.2 单调函数的平方积分下界不成立}
\noindent\auditref{p186-b04}{第 195 套试卷；原书第 185 页}\quad 结论错误\par
取 $f(x)=x/4-1$（$-1\le x\le0$）、$f(x)=x-1$（$0\le x\le1$）。连续严格递增，$f(0)=-1$，而积分为 $77/48<2f(0)^2=2$。
\par\textbf{修订建议：}不能仅凭单调性得到所写下界，需核对附加条件或系数。\par
\Needspace{8\baselineskip}
\subsection*{R041\quad 六题梯度定义的第二个下标笔误}
\noindent\auditref{p064-b18}{第 63 套试卷；原书第 63 页}\quad 题面异常\par
原图将两个梯度的第二分量分母均印为 $\partial g$。自变量为 $x,y$，应分别是 $\partial f/\partial y$ 和 $\partial g/\partial y$。
\par\textbf{修订建议：}两个分母中的 $g$ 应为 $y$；正文忠实保留原书。\par
\Needspace{8\baselineskip}
\subsection*{R042\quad 计算题8第一项重复微分 dz}
\noindent\auditref{p076-b06}{第 76 套试卷；原书第 75 页}\quad 题面异常\par
原图第一项确印为 $xy^3\,\mathrm dz\,\mathrm dz$，不是标准第二型曲面积分的微分对。其余两项依次为 $\mathrm dz\,\mathrm dx$ 与 $\mathrm dx\,\mathrm dy$。
\par\textbf{修订建议：}第一项很可能应为 $\mathrm dy\,\mathrm dz$，作为疑似笔误登记。\par
\Needspace{8\baselineskip}
\subsection*{R043\quad 四(1) 不一致连续判据未要求正常数}
\noindent\auditref{p075-b19}{第 75 套试卷；原书第 74 页}\quad 条件缺失\par
题面只写常数 $a$。若允许 $a=0$，取两列相同的点，则任何函数均满足所述序列条件，甚至恒零函数也满足，不能等价于不一致连续。
\par\textbf{修订建议：}应明确 $a>0$，并避免维数和数列指标共用同一符号造成歧义。\par
\Needspace{8\baselineskip}
\subsection*{R044\quad 第9题条件直接互相矛盾}
\noindent\auditref{p084-b16}{第 84 套试卷；原书第 83 页}\quad 题面异常\par
条件(1)写 $f(x)\ne0$，条件(2)写对 $n=0,1,\ldots$ 有 $f^{(n)}(x)=0$。取 $n=0$ 即得 $f(x)=0$，故不存在满足全部条件的函数。
\par\textbf{修订建议：}核实第二条是否应为某一端点的导数值；不可自行把 $x$ 改为 $0$。\par
\Needspace{8\baselineskip}
\subsection*{R045\quad 第7题三重积分上限未形成单纯形}
\noindent\auditref{p085-b10}{第 85 套试卷；原书第 84 页}\quad 结论错误\par
原书积分上限分别为 $x,y,z$，却声称结果只依赖于 $x$。取 $f\equiv1$，原式左侧为 $xyz$，右侧为 $x^3/6$，一般不相等。
\par\textbf{修订建议：}若原意为逐层积分，上限及微分变量应一并核实。\par
\Needspace{8\baselineskip}
\subsection*{R046\quad 第9(1)题 Gaussian 积分微分方程系数颠倒}
\noindent\auditref{p087-b13}{第 87 套试卷；原书第 86 页}\quad 结论错误\par
分部积分给出 $F'(x)=-2xF(x)$，且 $F(x)=\sqrt{\pi}e^{-x^2}/2$。代入原式得到 $-3xF(x)$，在 $x\ne0$ 不为零。
\par\textbf{修订建议：}应为 $F\,\prime(x)+2xF(x)=0$。\par
\Needspace{8\baselineskip}
\subsection*{R047\quad 第17题正项递增与级数收敛矛盾}
\noindent\auditref{p161-b19}{第 168 套试卷；原书第 160 页}\quad 题面异常\par
若 $a_{n+1}>a_n>0$，则 $a_n\ge a_1>0$，通项不趋零，级数不可能收敛。
\par\textbf{修订建议：}很可能应为正项递减；核对原始回忆后才能改入题面。\par
\Needspace{8\baselineskip}
\subsection*{R048\quad 第10题参数范围写成 0<m<m}
\noindent\auditref{p163-b11}{第 170 套试卷；原书第 162 页}\quad 题面异常\par
原图确印 $0<m<m$，严格不等式不可能成立，所求区域的参数范围没有有效定义。
\par\textbf{修订建议：}疑应为 $0<m<n$，正文保留原文并登记。\par
\Needspace{8\baselineskip}
\subsection*{R049\quad 七题任意实指数在负半轴未定义}
\noindent\auditref{p011-b21}{第 3 套试卷；原书第 10 页}\quad 条件缺失\par
题面允许任意 $\delta\ge0$，却在原点的二维邻域使用 $x^{1+\delta}$。如 $\delta=1/2$，负的 $x$ 上该实数幂无定义，不能直接讨论全邻域上的性质。
\par\textbf{修订建议：}核实是否漏绝对值或原有定义域限制；不同修订影响导数结论。\par
\Needspace{8\baselineskip}
\subsection*{R050\quad 第4题正弦四次幂缺自变量}
\noindent\auditref{p015-b02}{第 7 套试卷；原书第 14 页}\quad 题面异常\par
原图第一项确为 $\int\sin^4\,\mathrm dx$，没有写出正弦自变量。
\par\textbf{修订建议：}很可能漏 $x$；不凭常见题型静默补题。\par
\Needspace{8\baselineskip}
\subsection*{R051\quad 九(2) 平方根迭代并不一致收敛}
\noindent\auditref{p158-b21}{第 165 套试卷；原书第 157 页}\quad 结论错误\par
每个固定 $x\in(0,1)$ 都有 $f_n(x)=f_1(x)^{1/2^{n-1}}\to1$；但由 $f_1(x)<\sqrt{x}$ 可知对每个固定 $n$，$x\to0^+$ 时 $f_n(x)\to0$，故 $\sup_{0<x<1}|f_n(x)-1|=1$。
\par\textbf{修订建议：}应改为不一致收敛；在 $[\delta,1)$（$\delta>0$）上可得到一致收敛。\par
\Needspace{8\baselineskip}
\subsection*{R052\quad 第16题积分归一化因子错误}
\noindent\auditref{p170-b22}{第 179 套试卷；原书第 169 页}\quad 结论错误\par
固定 $a>b>0$，用积化和差及 Dirichlet 积分可得不带前因子的对称积分趋于 $\pi$；原式乘 $1/T$ 后趋于 $0$，而右侧趋于 $1$，互相矛盾。
\par\textbf{修订建议：}前因子很可能应为 $1/\pi$；应同时注明渐近过程为 $T\to+\infty$。\par
\subsection*{回忆缺失}
原文中“忘了”“具体数据忘了”等缺失说明均原样保留；完整索引见项目中的题目审阅记录。不可根据常见题型猜补数据。
